% Book 4, Year 2, chapter 20 -- The application layer: DNS, HTTP and TLS.
% Provenance (never printed): sources/book4/BRIEFS.md, chapter 20.
\chapter{The Application Layer: DNS, HTTP and TLS}\label{ch:b2:application-layer}

A program asks for the page at \code{http://www.brenmoor.test/}, the way a browser does when an address is typed:

\omoutput{code/bachelor-2/20-application-layer/outputs/visit.txt}{A first visit: a name lookup, a connection, a
request, a page.}\label{out:b2:application-layer:hook}

Three protocols took part. The \emph{name system} turned the name into an address, TCP opened a connection to it, and
HTTP carried the request and the page. All three servers are this chapter's own programs, running on this machine:
the name server listens on UDP port 44202 and the web server on TCP port 44203. This chapter builds them,
speaking the real formats, then adds the layer that makes the exchange private, TLS, and counts what a first visit
costs in round trips.

\begin{recall}
Book 1: a name such as \code{library.brenmoor.example} is looked up in the name system, which answers with an
address. Book 2: an HTTP request names a method and a path, and the response carries a status, headers and a body; a
secure connection uses keys and a certificate that a trusted authority has signed. \cref{ch:b2:sockets}: \omterm{def:b2:sockets:socket}{sockets} and
ports. \cref{ch:b2:transport-layer}: UDP datagrams, TCP's handshake and its byte stream.
\end{recall}

\section{The name system}\label{sec:b2:application-layer:dns}

% ledger: b2:test-tld
% ledger: b2:dns-zone
\begin{definition}[Domain name, name system, zone]\label{def:b2:application-layer:dns}
A \emph{domain name}\index{domain name} is a sequence of labels written from the most specific to the most general,
such as \code{www.brenmoor.test}; the names form a tree whose root has the empty name. The \emph{domain name system}
(DNS)\index{domain name system} is the distributed database that maps names to data such as addresses. It is cut into
\emph{zones}\index{zone (DNS)}: a zone is a part of the tree whose data one organisation keeps, and it may
\emph{delegate} a subtree to another zone.
\end{definition}

The name \code{.test} is reserved for testing, so no real server answers for it: the town's zone,
\code{brenmoor.test}, lives only in this chapter's programs.

\begin{definition}[Resolver, authoritative server]\label{def:b2:application-layer:resolver}
An \emph{authoritative server}\index{authoritative server} of a zone answers for the names of that zone, and for the
names of a delegated subtree it answers with a \emph{referral}: the names and addresses of that subtree's servers. A
\emph{resolver}\index{resolver} finds the data for a name. In \emph{iterative resolution} it asks the servers itself,
following referrals from the root down; in \emph{recursive resolution} it asks another resolver to do the whole job
and send only the final answer.
\end{definition}

\begin{omfigure}
\begin{tikzpicture}[font=\scriptsize, n/.style={draw, rounded corners=2pt, inner sep=2pt, fill=white},
    q/.style={-{Stealth[length=1.6mm]}, omThm, thick}]
\node[n] (root) at (4,3.3) {root (empty name)};
\node[n] (ex) at (2.4,2.2) {\code{example}}; \node[n] (test) at (5.2,2.2) {\code{test}};
\node[n] (bren) at (5.2,1.1) {\code{brenmoor.test}};
\node[n] (www) at (2.4,0) {\code{www}}; \node[n] (mail) at (4.0,0) {\code{mail}};
\node[n] (library) at (5.6,0) {\code{library}}; \node[n] (ns) at (7.2,0) {\code{ns}};
\draw (root) -- (test); \draw (root) -- (ex); \draw (test) -- (bren);
\foreach \l in {www, mail, library, ns} {\draw (bren) -- (\l);}
\draw[omDef, dashed, rounded corners=6pt] (1.7,-0.8) rectangle (7.9,1.45);
\node[text=omDef, anchor=south east] at (7.85,-0.77) {zone \code{brenmoor.test}};
\node[n, fill=omDef!12] (res) at (10,2.2) {resolver};
\draw[q] (res) to[bend right=12] node[above, text=omThm] {1} (root.east);
\draw[q] (res) -- node[above, text=omThm] {2} (test);
\draw[q] (res) to[bend left=10] node[above, text=omThm] {3} (bren.east);
\end{tikzpicture}
\par\smallskip{\small Part of the name tree. The root's servers refer a \omterm{def:b2:application-layer:resolver}{resolver} to those of \code{test} (query 1),
which refer it to those of \code{brenmoor.test} (query 2), whose zone holds the answer (query 3).}
\end{omfigure}

\begin{definition}[Resource record, record lifetime]\label{def:b2:application-layer:record}
% ledger: b2:dns-types
% ledger: b2:dns-aaaa
A \emph{resource record}\index{resource record} is one fact about a name: the name, a type, a \emph{lifetime} (TTL,
\omterm{def:b2:network-layer:ttl}{time to live}) and data. The common types are A (an IPv4 address, type 1), AAAA (an \omterm{def:b2:network-layer:ipv6}{IPv6 address}, 28), NS (a name
server of the zone, 2), CNAME (the name is an alias of another, 5), MX (a mail server, with a preference, 15) and
TXT (text, 16). A \omterm{def:b2:application-layer:resolver}{resolver} may keep a record in its cache for the record's lifetime, in seconds, and no longer.
\end{definition}

\begin{method}[Iterative resolution]\label{met:b2:application-layer:resolve}
To find the records of type $t$ for a name: start with a root server. Ask the current server for the name and $t$. If
the reply holds answers, return them (an alias, CNAME, comes with the records of its target when the zone holds them).
If it says the name does not exist, stop. If it is a referral, take one of the named servers, with its address from
the reply's additional records (the \emph{glue}), and ask it in turn.
\end{method}

\omoutput[1-1,9-23]{code/bachelor-2/20-application-layer/outputs/resolve.txt}{Iterative resolution on a toy
network of three servers, in the wire format of the next section: an alias, and a name that does not exist.}\label{out:b2:application-layer:resolve}

The servers of \code{zone.py} answer from a list of records; the referral comes from an NS record whose name lies
below the server's own zone. Each question costs one round trip to a server, so a \omterm{def:b2:application-layer:resolver}{resolver} that remembers nothing
pays three round trips for every name of the town.

\subsection{The message format}

% ledger: b2:dns-header
% ledger: b2:dns-port
% ledger: b2:dns-limits
Queries and answers share one format. A 12-byte header holds an identifier, which the answer copies so that the client
can match it with its question, sixteen bits of flags, and four counts; then come the question, and three lists of
records: the answers, the authority (NS records of a referral), and additional records (glue). A message travels in one
UDP datagram to port 53, or over TCP when it exceeds 512 bytes.

\begin{omfigure}
\begin{bytefield}[bitwidth=1.55em]{16}
\bitheader{0,1,4,5,6,7,8,9,11,12,15} \\
\bitbox{16}{\scriptsize identifier} \\
\bitbox{1}{\tiny QR} & \bitbox{4}{\scriptsize opcode} & \bitbox{1}{\tiny AA} & \bitbox{1}{\tiny TC} &
\bitbox{1}{\tiny RD} & \bitbox{1}{\tiny RA} & \bitbox{3}{\scriptsize Z} & \bitbox{4}{\scriptsize RCODE} \\
\bitbox{16}{\scriptsize QDCOUNT: questions} \\
\bitbox{16}{\scriptsize ANCOUNT: answers} \\
\bitbox{16}{\scriptsize NSCOUNT: authority records} \\
\bitbox{16}{\scriptsize ARCOUNT: additional records}
\end{bytefield}
\par\smallskip{\small The DNS header, six 16-bit fields. QR is 1 in an answer; AA says the server is authoritative;
TC that the message was truncated; RD that recursion is desired, RA that it is available; RCODE 3 means that the name
does not exist.}
\end{omfigure}

% ledger: b2:dns-compression
A name is written as its labels, each preceded by its length (at most 63), and ends with a zero byte, the empty label
of the root. Since the names of one answer repeat the same suffixes, a name may end with a \emph{pointer}: two bytes
whose first two bits are 1, the other 14 giving the offset, from the start of the message, of a suffix written
earlier.

\omcode{code/bachelor-2/20-application-layer/listings/dns_message.py}{10}{23}{Writing a name, with
compression}\label{lst:b2:application-layer:encode}

\omcode{code/bachelor-2/20-application-layer/listings/dns_message.py}{26}{41}{Reading a name: pointers must go
back}\label{lst:b2:application-layer:read}

A pointer could point anywhere, including to itself, and a careless reader would then loop for ever on a 14-byte
message from anyone. The reader above refuses any pointer that does not go back.

\begin{proposition}[Compression terminates]\label{prop:b2:application-layer:compression}
If every pointer points to an offset strictly before the start of the stretch of labels in which it appears, reading a
name from a message of $n$ bytes ends after at most $n$ steps.
\end{proposition}

\begin{proof}
Reading proceeds in stretches: a stretch starts at the name's position or at a pointer's target, and moves forward
label by label until a zero byte or a pointer. Within a stretch every step moves forward, and a stretch ends before
the end of the message. A pointer starts a new stretch at an offset strictly smaller than the start of the current
one, so the starts of the stretches strictly decrease. Every byte read lies at or after the start of its stretch and
before the start of the previous one (or the end of the message for the first stretch), so the stretches read
disjoint ranges of bytes, and the total number of steps is at most $n$.
\end{proof}

\omoutput{code/bachelor-2/20-application-layer/outputs/dns_message.txt}{A query and an answer, encoded and decoded
again.}\label{out:b2:application-layer:message}

With compression the answer takes 84 bytes instead of 143. On a real \omterm{def:b2:sockets:socket}{socket}, the zone's server answers queries for any
of its records; a name it does not hold gets RCODE 3, ``no such name'':

\omoutput{code/bachelor-2/20-application-layer/outputs/dns_demo.txt}{Five queries to the brenmoor.test server on
127.0.0.1, UDP port 44200.}\label{out:b2:application-layer:demo}

\subsection{Caching}

\omterm{def:b2:application-layer:resolver}{Resolvers} keep answers and referrals for their lifetimes. A machine asks one nearby \omterm{def:b2:application-layer:resolver}{resolver}, recursively; that
\omterm{def:b2:application-layer:resolver}{resolver} resolves iteratively, and its cache serves every machine behind it. The gain depends on how often a name is
asked for: if queries for a name arrive at random at rate $r$ per second and its record lives $T$ seconds, the fraction
answered from the cache is $rT/(1 + rT)$ (\cref{exo:b2:application-layer:8}).

\omoutput{code/bachelor-2/20-application-layer/outputs/resolver_cache.txt}{The fraction of queries a cache answers,
simulated over 200\,000 queries.}\label{out:b2:application-layer:cache}

A name asked once an hour with a lifetime of an hour still saves half the queries; a busy name is almost never asked
upstream. The price is staleness: a record changed at its zone may be served from caches for up to its old lifetime,
which is why \omterm{def:b2:application-layer:dns}{zones} lower the lifetime of a record some time before changing it.

\section{HTTP}\label{sec:b2:application-layer:http}

% ledger: b2:http-ports
An HTTP request is lines of text: a method, a path and a version, then headers, then an empty line, and a body for
methods that send one. The response starts with a status line, then headers and the body. The default port is 80, and
443 for HTTP over TLS. A request by hand, then two ways of fetching twenty pages:

\omoutput{code/bachelor-2/20-application-layer/outputs/http_demo.txt}{A request typed by hand on a \omterm{def:b2:sockets:socket}{socket}, then twenty
requests through two client libraries.}\label{out:b2:application-layer:http}

\begin{definition}[Persistent connection, head-of-line blocking]\label{def:b2:application-layer:persistent}
% ledger: b2:http-persistent
A \emph{persistent connection}\index{persistent connection} carries several requests and their responses, one after
the other, instead of one; it is HTTP/1.1's default, and a \code{Connection: close} header announces that the
connection will be closed after the response. \emph{Head-of-line blocking}\index{head-of-line blocking} is the delay a
response suffers because something before it on the same connection, or in the same ordered stream, has not finished.
\end{definition}

On a \omterm{def:b2:application-layer:persistent}{persistent connection} the client must know where each body ends, so the server sends either a
\code{Content-Length} header or, when it does not know the length in advance, a \emph{chunked} body: chunks, each
preceded by its size in hexadecimal, and a last chunk of size 0 (\cref{exo:b2:application-layer:5}). The town's
server reads requests one after the other on each connection:

\omcode{code/bachelor-2/20-application-layer/listings/http_server.py}{11}{34}{One connection of the web server, from
the first request to the close}\label{lst:b2:application-layer:server}

\code{urllib} asked for a close each time, so it paid a TCP handshake per page; \code{http.client} kept one
connection for all twenty. On a connection of HTTP/1.1, though, the second response cannot start before the first is
complete. Browsers open several connections at once to fetch a page's parts in parallel.

\begin{definition}[Stream multiplexing]\label{def:b2:application-layer:streams}
% ledger: b2:http2-frame
In HTTP/2, a connection carries numbered \emph{streams}\index{stream (HTTP/2)}, one per request and its response,
whose data travel in \emph{frames}: each frame starts with a 9-byte header giving its length, its type (such as DATA 0
and HEADERS 1), flags and the number of its stream. Frames of different streams may be interleaved, so responses
progress together on one connection: this is \emph{stream multiplexing}\index{stream multiplexing}.
\end{definition}

\begin{omfigure}
\begin{bytefield}[bitwidth=0.8em]{40}
\bitheader{0,23,24,31,32,39} \\
\bitbox{24}{\scriptsize length of the payload} & \bitbox{8}{\scriptsize type} & \bitbox{8}{\scriptsize flags} \\
\bitbox{1}{\tiny R} & \bitbox{31}{\scriptsize stream identifier} & \bitbox[]{8}{}
\end{bytefield}
\par\smallskip{\small The 9-byte header of an HTTP/2 frame (72 bits, drawn on two rows). The payload, by default at
most $2^{14}$ bytes, follows.}
\end{omfigure}

\omterm{def:b2:transport-layer:demux}{Multiplexing} removes \omterm{def:b2:application-layer:persistent}{head-of-line blocking} between responses, but not below them: the frames still travel in one TCP
byte stream, and a lost TCP segment holds back everything after it until it is sent again. A model of four pages of 8
packets each, with a round trip of 50 ms, and the fifth packet lost:

\omoutput{code/bachelor-2/20-application-layer/outputs/hol_blocking.txt}{When each page is complete, in a model with
one lost packet (\code{hol_blocking.py}).}\label{out:b2:application-layer:hol}

HTTP/1.1 fetches the pages in turn, and the loss delays all of them. HTTP/2 brings them in parallel, but the loss
delays all four to 107 ms, when the lost packet's copy arrives. With HTTP/3, each stream is ordered on its own, and
only the page that lost a packet waits (\cref{exo:b2:application-layer:10}).

\begin{dated}{2026}
% ledger: b2:http3
HTTP/3, standardised in June 2022, carries HTTP over QUIC (\cref{ch:b2:transport-layer}): its \omterm{def:b2:application-layer:streams}{streams} are QUIC
\omterm{def:b2:application-layer:streams}{streams}, each delivered in order on its own, so that a loss on one stream does not hold back the others. It is used
alongside HTTP/2 and HTTP/1.1, which remain in wide use; a server announces that it speaks HTTP/3 in its HTTP/1.1 or
HTTP/2 responses.
\end{dated}

\section{Electronic mail}\label{sec:b2:application-layer:mail}

% ledger: b2:smtp
Mail moves between servers by SMTP, over TCP port 25, in a dialogue of text lines. To deliver a letter to
\code{clerk@example.net}, a server looks up the MX records of \code{example.net}, tries the one with the lowest
preference value first, connects to it, and the two exchange:

\begin{center}\small
\begin{tabular}{@{}l l l@{}}
\textbf{client} & \textbf{server} & \textbf{meaning} \\
& \code{220 mail.example.net} & ready \\
\code{EHLO post.example.org} & \code{250} & greeting \\
\code{MAIL FROM:<ada@example.org>} & \code{250} & the sender \\
\code{RCPT TO:<clerk@example.net>} & \code{250} & one recipient (several allowed) \\
\code{DATA} & \code{354} & send the letter \\
the letter, then a line holding only \code{.} & \code{250} & accepted \\
\code{QUIT} & \code{221} & closing \\
\end{tabular}
\end{center}

The three-digit codes are for programs, the text after them for people; the first digit says whether the command
succeeded (2), needs more (3), or failed (4 temporarily, 5 for good). The letter itself, with its \code{From},
\code{To} and \code{Subject} lines, is only data to SMTP: nothing in the protocol checks that the \code{From} line is
true, which is why mail systems added signatures and checks of their own.

\section{TLS}\label{sec:b2:application-layer:tls}

Everything so far crossed the network in clear: anyone on the path could read the page or change it. TLS adds a
layer between TCP and the application.

\begin{definition}[Transport Layer Security]\label{def:b2:application-layer:tls}
% ledger: b2:tls13-record
\emph{Transport Layer Security} (TLS)\index{Transport Layer Security} is a layer that runs over a reliable byte stream
such as TCP and gives the application above it a byte stream that is private (encrypted), protected against
modification, and attached to a verified name. It cuts the application's bytes into \emph{TLS
records}\index{TLS record}, each of at most $2^{14}$ bytes of data, sent encrypted after a 5-byte header (a content
type, a version and a length).
\end{definition}

% ledger: b2:tls13-handshake
Before any record of data, a \emph{handshake} agrees on keys and proves the server's identity. In TLS 1.3 it takes one
round trip:

\begin{omfigure}
\begin{tikzpicture}[font=\scriptsize, w/.style={-{Stealth[length=1.8mm]}, thick}]
\node[font=\small\bfseries] at (0,0.3) {client}; \node[font=\small\bfseries] at (9,0.3) {server};
\draw[gray] (0,0) -- (0,-4.2); \draw[gray] (9,0) -- (9,-4.2);
\draw[w, omDef] (0,-0.3) -- node[above, sloped] {ClientHello (versions, ciphers, a key share)} (9,-0.8);
\draw[w, omDef] (9,-1.6) -- node[above, sloped, align=center] {ServerHello (a key share), then encrypted:\\
EncryptedExtensions, Certificate, CertificateVerify, Finished} (0,-2.1);
\draw[w, omDef] (0,-2.5) -- node[above, sloped] {Finished, then the HTTP request (encrypted)} (9,-3.0);
\draw[w, omThm] (9,-3.4) -- node[above, sloped] {the response (encrypted)} (0,-3.9);
\node[anchor=east, align=right] at (-0.1,-2.1) {keys known,\\ server verified};
\end{tikzpicture}
\par\smallskip{\small The TLS 1.3 handshake on an open \omterm{def:b2:transport-layer:tcp}{TCP connection}. After the first round trip both sides have the
keys, and the client has checked the server's certificate and its signature; the request goes with the client's
Finished.}
\end{omfigure}

The certificate binds the server's public key to its name, and is signed by an authority that the client trusts; the
mathematics of keys and signatures are the Year 3 volume's. Python's \code{ssl} module provides TLS over a \omterm{def:b2:sockets:socket}{socket}. The
session below makes a certificate for \code{www.brenmoor.test} with \code{openssl} (a self-signed one, which the
client is told to trust), then runs an echo pair and a client that asks for another name:

\omoutput[1-7]{code/bachelor-2/20-application-layer/outputs/tls_demo.txt}{TLS on 127.0.0.1: a verified connection,
then a certificate presented for the wrong name.}\label{out:b2:application-layer:tls}

The second client refuses the connection: the certificate is valid, but not for \code{mail.brenmoor.test}. Checking
the name is as important as checking the signature; a client that skipped it would accept any server holding any
certificate.

\begin{omfigure}
\begin{tikzpicture}[font=\scriptsize, b/.style={draw, minimum height=8mm, inner sep=2pt, align=center}]
\node[b, minimum width=16mm, anchor=west] (ip) at (0,0) {IP header\\ addresses};
\node[b, minimum width=16mm, anchor=west] (tcp) at (ip.east) {TCP header\\ ports};
\node[b, minimum width=16mm, anchor=west] (tls) at (tcp.east) {TLS record\\ header};
\node[b, minimum width=50mm, anchor=west, fill=omDef!20] (data) at (tls.east) {HTTP request or response,\\
content type, authentication tag};
\draw[decorate, decoration={brace, amplitude=4pt, mirror}] (ip.south west) -- node[below=4pt] {seen by every router
on the path} (tls.south east);
\draw[decorate, decoration={brace, amplitude=4pt}] (data.north west) -- node[above=4pt] {encrypted: only the two
ends read it} (data.north east);
\end{tikzpicture}
\par\smallskip{\small A packet of a TLS connection. TLS hides the content and protects it against change, but not who
talks to whom, when, or how much.}
\end{omfigure}

What TLS does not hide matters too. The addresses and ports are visible to every \omterm{def:b2:network-layer:routing}{router}, and so are the sizes and
times of the records. The name lookup before the connection normally travels in clear, and the ClientHello
usually carries the name of the site in clear, so that a server holding several sites can pick the right certificate.

\subsection{The cost of a first visit}

Every protocol of this chapter waits for an answer before the next can start, so a first visit is a chain of round
trips.

\begin{proposition}[Round trips of a visit]\label{prop:b2:application-layer:round-trips}
Let a page be small enough to arrive in one round trip, and let RTT be the round-trip time to the server. A first
visit over HTTPS, with TCP and TLS 1.3, takes the time of the name lookup plus $3\,\mathrm{RTT}$. A visit on a new
connection with the name cached on the machine takes $3\,\mathrm{RTT}$, and a request on a connection already open
takes $1\,\mathrm{RTT}$.
\end{proposition}

\begin{proof}
The client cannot connect before it knows the address. TCP's handshake costs one round trip: the client can send only
after the SYN+ACK arrives (its last ACK travels with the ClientHello). The client can send the request only after the
server's Finished arrives, one more round trip. The response comes back one round trip after the request leaves. Each
of these waits is a full round trip and none can overlap the next, which gives $3\,\mathrm{RTT}$; on an open
connection only the last remains.
\end{proof}

\omoutput{code/bachelor-2/20-application-layer/outputs/budget.txt}{The budget with a round trip of 80 ms to the
server, 10 ms to the \omterm{def:b2:application-layer:resolver}{resolver}, and 40 ms from the \omterm{def:b2:application-layer:resolver}{resolver} to each name server.}\label{out:b2:application-layer:budget}

A cold first visit spends 130 ms of its 370 on the name, which is why \omterm{def:b2:application-layer:resolver}{resolvers} cache, and the rest on handshakes,
which is why browsers keep connections open and why TLS 1.3 cut the handshake to one round trip.

\begin{inpractice}
Opening a page usually costs more than this: a page names dozens of other resources (pictures, style sheets,
scripts), often on other servers, each with its own name to look up and its own connection. Shortening the chain of
round trips (\omterm{def:b2:application-layer:persistent}{persistent connections}, caches near the users, fewer servers per page) often saves more time than a faster
link does.
\end{inpractice}

\omimg{book4/ai/b2-application-layer-datacentre-aisle.jpg}{0.6\linewidth}{An aisle of servers in a data centre. Busy
sites run many copies of their servers in several such buildings, and the name system can send each user to a nearby
one. AI illustration.}

\section{Exercises}\label{sec:b2:application-layer:exercises}

\begin{exercise}[$\star$]\label{exo:b2:application-layer:1}
Decode the DNS header \code{2a 17 85 83 00 01 00 00 00 00 00 00}: is it a question or an answer, what flags are set,
and what does it say about the name asked?
\end{exercise}

\begin{exercise}[$\star$]\label{exo:b2:application-layer:2}
A message contains, from offset 12, the bytes \code{03 77 77 77 c0 0c}. What does a reader that follows every pointer
do with it? Bound the number of pointers a reader must follow in a message of $n$ bytes, with or without the rule of
\cref{prop:b2:application-layer:compression}.
\end{exercise}

\begin{exercise}[$\star$]\label{exo:b2:application-layer:3}
Write, byte by byte, a query with identifier \code{0b0e} that asks for the AAAA records of \code{www.brenmoor.test},
with recursion desired.
\end{exercise}

\begin{exercise}[$\star$]\label{exo:b2:application-layer:4}
A machine asks its \omterm{def:b2:application-layer:resolver}{resolver} for \code{mail.brenmoor.test}. How many queries does the \omterm{def:b2:application-layer:resolver}{resolver} send if its cache is
empty, if it holds the referral to \code{brenmoor.test}, and if it holds the answer? How many round trips does the
machine wait in each case?
\end{exercise}

\begin{exercise}[$\star\star$]\label{exo:b2:application-layer:5}
Write \code{dechunk(data)}, which decodes a chunked body and returns it with the number of bytes used, and raises an
error if the coding is cut short.
\end{exercise}

\begin{exercise}[$\star\star$]\label{exo:b2:application-layer:6}
This server announces 12 bytes and sends 10. What does a client do, and why?
\omcode{code/bachelor-2/20-application-layer/listings/exo_content_length.py}{9}{13}{A wrong length}
\end{exercise}

\begin{exercise}[$\star\star$]\label{exo:b2:application-layer:7}
A first version of the server answered one request per connection and knew only \code{GET}. Write the changes that
keep the connection open unless the client asks to close it or speaks HTTP/1.0, and that answer \code{HEAD}, which
sends the headers of \code{GET} without the body.
\end{exercise}

\begin{exercise}[$\star\star$]\label{exo:b2:application-layer:8}
Queries for a name arrive as a Poisson process of rate $r$, and a \omterm{def:b2:application-layer:resolver}{resolver} caches the answer for $T$ seconds after
each query it sends upstream. Show that the fraction of queries answered from the cache is $rT/(1 + rT)$.
\end{exercise}

\begin{exercise}[$\star\star$]\label{exo:b2:application-layer:9}
Write a \omterm{def:b2:application-layer:resolver}{resolver} that resolves iteratively, keeps answers and referrals for their lifetimes, and starts each
resolution from the closest zone whose servers it knows.
\end{exercise}

\begin{exercise}[$\star\star\star$]\label{exo:b2:application-layer:10}
In the model of \cref{out:b2:application-layer:hol}, packets leave one per millisecond from 25 ms, and a lost packet's
copy arrives $3 \times 25$ ms after the third packet behind it leaves. Compute when each page is complete with HTTP/2
and with HTTP/3, and explain why HTTP/1.1 is slowest.
\end{exercise}

\begin{exercise}[$\star\star\star$]\label{exo:b2:application-layer:11}
Write a TLS echo pair on 127.0.0.1: the server uses a certificate and key from files, and the client trusts only that
certificate and checks the server's name.
\end{exercise}

\begin{exercise}[$\star\star\star$]\label{exo:b2:application-layer:12}
Measure, after the handshake, how many bytes TLS sends for messages of 1 byte to 100\,000 bytes, and explain the
extra bytes per record. What fraction of the traffic is overhead for messages of 1, 100 and 16\,384 bytes?
\end{exercise}

\section{Problem: Brenmoor goes online}\label{sec:b2:application-layer:problem}

\begin{problem}[{Weekend problem --- Brenmoor goes online}]\label{pb:b2:application-layer:1}
The town puts its name server, its \omterm{def:b2:application-layer:resolver}{resolver} and its web site online, and wants to know what a first visit costs.

\medskip\noindent\textbf{Part I --- The zone and its server.}\par\nopagebreak
\begin{enumerate}
\item List the records of \code{zone.py}'s town zone by type, and say what each is for.
\item Why must the zone \code{test} hold the address of \code{ns.brenmoor.test}, although the name lies in
\code{brenmoor.test}?
\item In \cref{out:b2:application-layer:demo}, why is the answer for the MX query longer than the query by only 21
bytes?
\item A record of \code{library} is a CNAME. What does the server send for its A record, and why in one answer?
\item What does RCODE 3 tell a \omterm{def:b2:application-layer:resolver}{resolver} that an empty answer with RCODE 0 does not?
\item Why does the server reject a name whose pointer does not go back, instead of following it?
\end{enumerate}

\medskip\noindent\textbf{Part II --- A resolver with a cache.}\par\nopagebreak
\begin{enumerate}[resume, beginpenalty=10000]
\item How many queries does an empty \omterm{def:b2:application-layer:resolver}{resolver} send for \code{www}, then for \code{mail}, then for \code{www} again
within the hour?
\item After three days without queries, how many does it send for \code{www}?
\item The town's site is asked for 1000 times a day. With a lifetime of 3600 s, what fraction of the lookups reach
the town's server?
\item The town moves its web server. How long may some users still reach the old address, and how can the town
shorten it?
\item Why does a \omterm{def:b2:application-layer:resolver}{resolver} use the shortest lifetime of a set of records for the whole set?
\item Why do machines ask a \omterm{def:b2:application-layer:resolver}{resolver} recursively, rather than resolve iteratively themselves?
\end{enumerate}

\medskip\noindent\textbf{Part III --- A web server with persistent connections.}\par\nopagebreak
\begin{enumerate}[resume, beginpenalty=10000]
\item In \cref{out:b2:application-layer:http}, how many TCP handshakes did each client pay?
\item With a round trip of 80 ms, how much time did the \omterm{def:b2:application-layer:persistent}{persistent connection} save over twenty pages?
\item How does the server know where a request ends, and the client where a response ends?
\item Why can the server not answer the second request on a connection before the first?
\item In \cref{out:b2:application-layer:hol}, which pages does the loss delay with HTTP/2, and by how much?
\item Why does HTTP/3 need a transport other than TCP to avoid that delay?
\end{enumerate}

\medskip\noindent\textbf{Part IV --- The cost of a first page load.}\par\nopagebreak
\begin{enumerate}[resume, beginpenalty=10000]
\item Using \cref{prop:b2:application-layer:round-trips}, how long does a cold first visit take in
\cref{out:b2:application-layer:budget}, and how much of it is the name?
\item How long does a revisit on a new connection take, and one on an open connection?
\item How many round trips would a TLS handshake of two round trips add to every new connection?
\item A page needs 30 resources from 3 servers. With one connection per server kept open, how many handshakes does it
pay?
\item What does TLS hide on such a visit, and what can an observer of the network still learn?
\item State the result: the round-trip budget of a first visit (name lookup, connection, handshake, request) and of a
revisit.
\end{enumerate}
\end{problem}
